En el presente informe se expone el análisis de dos problemas de la computación muy populares; la multiplicación de matrices y el ordenamiento de arreglos unidimensionales. Buscamos 
comparar la complejidad asintótica teórica de los diversos algoritmos, para los ordenamientos analizaremos \textsc{Merge Sort} \cite{gfg_merge_sort}, \textsc{Quick Sort} \cite{gfg_quick_sort}, \textsc{Patience Sort} \cite{gfg_patience_sorting} y \texttt{std::sort},
mientras que para la multiplicación de matrices veremos los algoritmos de \textsc{Naive} y \textsc{Strassen} \cite{medium_strassen}.
Evaluaremos el comportamiento de ambos problemas ante altos volúmenes de datos a través de distintas distribuciones, tamaños de entrada y dominios. Con esto buscamos evidenciar
cómo los factores del hardware impactan en los tiempos de ejecución, demostrando en la práctica las limitaciones de depender exclusivamente de la notación Big O
